\section{JAX-native SMC: once-per-run compilation} \label{sec:jax-native}

This section covers what is, in the authors' view, the framework's
single most consequential engineering decision: the JAX-native
re-implementation of the tempered-SMC outer loop in
\path{smc2fc/core/jax_native_smc.py}, eliminating the per-window JIT
recompilation cost of the BlackJAX-wrapped baseline and reducing the
end-to-end runtime of a multi-window inference run by roughly an
order of magnitude.

It also documents a failed experiment --- replacing HMC with NUTS
(\path{smc2fc/core/jax_native_smc_nuts.py}) --- so that future
contributors do not repeat it.

\subsection{The compile-cost problem with the BlackJAX-wrapped backend}

The reference implementation \path{smc2fc/core/tempered_smc.py} is a
thin wrapper around BlackJAX's
\texttt{blackjax.smc.tempered.build\_kernel}. BlackJAX bakes the
log-likelihood function into the closure of the returned kernel.
That is fine for a single run; the issue is that in the rolling-window
setting (Section~\ref{sec:mpc}) the log-likelihood closes over per-
window observation data, so a new kernel is built per window. Every
fresh kernel triggers a new \texttt{jax.jit} cache entry, and every
new cache entry costs a full HLO recompile --- on the model classes
the framework targets, in the order of $\sim 15$ seconds per window.

For a $27$-window benchmark that recompilation dominates the total
wall-clock; the actual sampling work is only a small fraction of the
visible cost. This is a structural problem of the BlackJAX API, not
of BlackJAX's HMC kernel itself --- the kernel is fine; the issue is
that it is rebuilt every window.

\subsection{The fix}

The fix has three parts, all inside
\path{smc2fc/core/jax_native_smc.py}:

\paragraph{(a) Module-level HMC kernel.} The HMC kernel itself is
built once at module import:
\begin{lstlisting}
HMC_KERNEL = hmc.build_kernel()  # smc2fc/core/jax_native_smc.py:53
HMC_INIT   = hmc.init
\end{lstlisting}
These are module-level constants, not per-call closures. The kernel's
trace cache key depends only on the shape and dtype of its arguments,
not on their values. Same shape/dtype across windows $\Rightarrow$ one
compile per process.

\paragraph{(b) Pytree-Partial wrapping of per-window callables.} The
per-window log-density and the prior log-density are passed in as
\texttt{jax.tree\_util.Partial}-wrapped callables. Partial is a
JAX-recognised pytree, so its bound-arg \emph{shapes/dtypes} enter the
trace cache key but its bound-arg \emph{values} do not. Across
windows the data values change but the data shapes don't, so the
cache hits (\path{jax_native_smc.py:121--130}).

\paragraph{(c) On-device tempering loop via \texttt{lax.while\_loop}.}
The full adaptive-tempering loop, including the ESS bisection that
chooses each $\delta\lambda$, is implemented inside a single
\texttt{lax.while\_loop} (\path{jax_native_smc.py:252}). The
ESS bisection is itself a fixed-iteration \texttt{lax.scan}
(\path{jax_native_smc.py:_solve_delta_for_ess}, 30 bisection steps).
There is no Python-side iteration, no Python-side branching, no
Python-side ESS solver. The whole thing is one XLA graph.

The combined effect: the first window's call pays the JIT compile
cost ($\sim 1.6$ seconds for the cached HMC kernel); every subsequent
window reuses the compiled artefact and pays only the sampling cost
($\sim 1.1$ seconds per window in the benchmark).

\subsection{Empirical speedup}

On a synthetic benchmark target with $256$ parameter particles and
$20$ tempering steps:
\begin{center}
\begin{tabular}{lr}
\toprule
Stage & HMC backend (jax-native) \\
\midrule
JIT compile + first run        & 1.63\,s \\
Cached run (subsequent windows) & 1.09\,s per window \\
\bottomrule
\end{tabular}
\end{center}

For the same benchmark, the BlackJAX-wrapped backend pays roughly
$15$ seconds of recompilation per window (figures will vary by
hardware and HLO cache state); the JAX-native backend amortises this
to a one-shot $1.6$\,s and is then bounded only by the cost of
sampling. The net speedup at $27$ windows is approximately one order
of magnitude.

\subsection{Failed experiment: NUTS as a drop-in replacement} \label{sec:nuts-failed}

The textbook upgrade from HMC is the No-U-Turn Sampler (NUTS), which
adapts the trajectory length per step based on a U-turn diagnostic. A
JAX-native NUTS variant of \path{jax_native_smc.py} was implemented
in \path{smc2fc/core/jax_native_smc_nuts.py} (only the kernel
differs --- everything else is structurally identical). Both kernels
were benchmarked side-by-side on the same synthetic target inside the
same conda environment.

\begin{center}
\begin{tabular}{lrr}
\toprule
Stage & HMC & NUTS \\
\midrule
JIT compile + first run        & 1.63\,s & 102.03\,s \\
Cached run (per window)         & 1.09\,s & 94.23\,s \\
\bottomrule
\end{tabular}
\end{center}

The NUTS version is approximately \emph{60$\times$ slower at compile
time and $85\times$ slower per cached window}.

The cause is structural, not parametric. NUTS dynamically determines
the trajectory length using an internal \texttt{while\_loop} (the
U-turn check). When this dynamic-control-flow loop is placed inside
a \texttt{jax.vmap} (over the $256$ parameter particles) and an outer
\texttt{lax.while\_loop} (the tempering schedule), JAX is forced to
either unroll the dynamic loop --- which produces an enormous XLA
graph --- or compile a deeply nested control-flow structure that the
GPU executes inefficiently. Either way the cost compounds across the
two outer loops.

The trajectory-length expansion is also unbounded in the worst case.
``$5$ NUTS steps per particle per tempering level'' (matching the HMC
default \texttt{num\_mcmc\_steps = 5}) translates into $5 \times 8 =
40$ gradient evaluations under HMC, but anywhere from $5 \times 1 =
5$ to $5 \times 1024 = 5120$ evaluations under NUTS depending on the
maximum tree depth.

\paragraph{Conclusion.} HMC remains the production rejuvenation
kernel. The NUTS implementation file
\path{smc2fc/core/jax_native_smc_nuts.py} is preserved as a documented
dead-end --- not deleted, so the next contributor who has the same
``why don't we just use NUTS?'' idea can read this section, see the
benchmark, and not repeat the experiment.

\subsection{How to switch backends}

In user code, the backends are interchangeable at the import line:

\begin{lstlisting}
# Production: JAX-native, compile-once, HMC.
from smc2fc.core.jax_native_smc import (
    run_smc_window_native, run_smc_window_bridge_native,
)

# Reference: BlackJAX wrapper, slower, easier to read.
from smc2fc.core.tempered_smc import (
    run_smc_window, run_smc_window_bridge,
)

# Documented dead-end (do not use; see Section ref{sec:nuts-failed}).
# from smc2fc.core.jax_native_smc_nuts import (
#     run_smc_window_native, run_smc_window_bridge_native,
# )
\end{lstlisting}

\subsection{Implementation pointers}

\begin{itemize}
  \item \path{smc2fc/core/jax_native_smc.py} ($\sim 430$ lines) ---
        production HMC backend.
  \item \path{smc2fc/core/jax_native_smc_nuts.py} ($\sim 430$ lines)
        --- the dead-end NUTS variant. Read the file's header
        docstring before considering it.
  \item \path{smc2fc/core/mass_matrix.py} ($20$ lines) --- the diagonal
        mass-matrix re-estimation discussed in Section~\ref{sec:mass-matrix}.
  \item \path{smc2fc/core/tempered_smc.py} ($\sim 370$ lines) ---
        BlackJAX-wrapped reference backend.
\end{itemize}

\subsection{Practical takeaway}

For any new application the framework targets, default to
\texttt{jax\_native\_smc.py}. The BlackJAX wrapper is useful only as
a reference when debugging an unfamiliar log-density (its slower
trace cycle gives you cleaner stack traces). The NUTS variant should
not be used.
